\documentclass[11pt,twoside]{amsart}
\usepackage{amssymb, amsmath, enumerate, palatino, hyperref}
\usepackage[normalem]{ulem}
\usepackage{fullpage}
\usepackage[T1]{fontenc}
\renewcommand{\labelitemi}{\guillemotright}
\usepackage{mathrsfs}

\usepackage{tikz}
\tikzstyle{ball} = [circle,shading=ball, ball color=black,
    minimum size=1mm,inner sep=1.3pt]

\theoremstyle{plain}
\newtheorem{prop}{Proposition}%[section]
\newtheorem{lemma}[prop]{Lemma}
\newtheorem{thm}[prop]{Theorem}
\newtheorem{obs}[prop]{Observation}
\newtheorem{app}[prop]{Application}
\newtheorem*{MainThm}{Main Theorem}
\newtheorem{cor}[prop]{Corollary}
\newtheorem{conj}[prop]{Conjecture}
\theoremstyle{remark}
\newtheorem{rmk}[prop]{Remark}
\newtheorem{prob}{Problem}
\newtheorem{bonus}[prop]{Bonus Problem}
\newtheorem{exc}{Exercise}
\theoremstyle{definition}
\newtheorem{ex}[prop]{Example}
\theoremstyle{definition}
\newtheorem{defn}[prop]{Definition}

\newcommand{\RR}{\mathbb{R}}
\newcommand{\ZZ}{\mathbb{Z}}
\newcommand{\CC}{\mathbb{C}}
\newcommand{\NN}{\mathbb{N}}
\newcommand{\QQ}{\mathbb{Q}}
\newcommand{\PP}{\mathbb{P}}

\newcommand{\cC}{\mathscr{C}}
\newcommand{\Ob}{\operatorname{Ob}}
\newcommand{\Mor}{\operatorname{Mor}}

\newcommand{\Set}{\operatorname{Set}}
\newcommand{\FinSet}{\operatorname{FinSet}}
\newcommand{\Vect}{\operatorname{Vect}}
\newcommand{\FinVect}{\operatorname{FinVect}}
\newcommand{\Mat}{\operatorname{Mat}}
\newcommand{\Gp}{\operatorname{Gp}}
\newcommand{\FinGp}{\operatorname{FinGp}}
\newcommand{\AbGp}{\operatorname{AbGp}}
\newcommand{\Ring}{\operatorname{Ring}}
\newcommand{\CommRing}{\operatorname{CommRing}}
\newcommand{\Field}{\operatorname{Field}}
\newcommand{\Top}{\operatorname{Top}}
\newcommand{\Aut}{\operatorname{Aut}}
\newcommand{\id}{\operatorname{id}}
\DeclareRobustCommand{\stir}{\genfrac\{\}{0pt}{}}

% For solutions.
\newenvironment{sol}{
  \medskip

  \noindent\textsc{Solution:}
}
{
  \medskip

}


\newcommand{\defeq}{:=}

\title{Math 113: Discrete Structures\\ Homework 35}
%\author{} %Remove the leading % and put your name here if using this .tex file as a template for your homework.

\begin{document}
\maketitle



\begin{prob}
  Find, with proof, the remainder of~$3^{1260126012601260129}$ upon division
  by~$14$.
\end{prob}

\begin{prob}\
  \begin{enumerate}[(a)]
    \item For~$n=10,11,$ and~$12$. List the fractions
      \[
	\frac{1}{n},\frac{2}{n},\dots,\frac{n}{n}
      \]
      after reducing each to lowest terms (canceling common factors in the
      numerator and denominator).
    \item Let $n$ be a natural number and consider the quantity
\[
  \psi(n)=\sum_{d\mid n}\varphi(d)
\]
which is the sum of the values $\varphi(d)$ where $d$ ranges through all the
positive divisors of $n$.  What is $\psi(n)$?  (Experiment, formulate a
conjecture, and prove it.) Your solution should consist of a precise statement
of your conjecture and a proof. The proof does not need to be elaborate.  It can
just be a statement of the general relevant phenomenon you observe in part~(a).
  \end{enumerate}
\end{prob}

\begin{prob} There are integers~$n$ such that~$-1$ has a square root in~$\ZZ/n\ZZ$.  To test this
  out, for each~$n\in\left\{ 2,3,\dots,13 \right\}$, find all solutions
  $x\in\left\{ 0,1,\dots,n-1 \right\}$ to the equation
  \[
    x^2\equiv -1\pmod{n}.
  \]
  You do not need to show your work.  It may help to note that~$-1\equiv n-1\pmod{n}$.
\end{prob}



\end{document}
